% Introduction
% Structure: Problem/Motivation -> Background -> Research Questions -> Contributions -> Paper Outline

The main concept that we will be investigating throughout this paper is speech. The way that speech is expressed and restricted is a significant factor into determining human behavior. It is equally important in understanding the health of public squares, whether they even exist, and who contributes to them. In this paper, we will first analyze and discuss this concept as a whole, and understand the correlation between free speech laws and the sharing of information. Once we have established this link, we will analyze whether or not a sense of surveillance has a chilling effect on the contribution of the uncensored sharing of information. At this point, we will look to Wikipedia for our data analysis, as it is the most commonly used example of a public square. Due to its decentralized, open nature, it is not linked to one government or entity that could take punitive measures for free speech. This model is also responsible for how truth is found on Wikipedia, with multiple potentially anonymous editors attempting to extract data from opinion, and come to an unbiased conclusion on different topics. This makes it useful for our analysis, as it is an online public square, one where contributions are documented. We will use Wikipedia as a microcosm for a public square, one with global users and try to understand how different freedom of speech laws affect contributions to these public squares. We do this by tracking where Wikipedia contributions are coming from via IP address, and breaking down contributions by country. We then take data from the V-Dem democracy index to quantify freedom of speech laws country by country. 

% Research questions
This paper investigates the following research questions:
\begin{enumerate}
    \item Does surveillance have a chilling effect on speech in public forums? Is this dependent on enforcement or just the feeling of being surveilled? How does this relate to Wikipedia as a case study?
    \item Do freedom of speech laws have a chilling effect on the number of anonymous Wikipedia contributions?
    \item How do governance and administrative controls impact trust and anonymity?
\end{enumerate}

% Contributions
The main contributions of this work are:
\begin{itemize}
    \item Providing quantitative evidence to link free speech laws and Wikipedia edits.
    \item Show how governance mechanisms can impact participation and trust.
    \item Provide insights into balancing privacy and protection against vandalism in online public squares.
\end{itemize}

% Paper outline
The rest of this paper is organized as follows. 
Section~\ref{sec:related} reviews related work on Wikipedia governance. 
Section~\ref{sec:methodology} describes our data and methods.
Section~\ref{sec:results} presents our findings.
Section~\ref{sec:discussion} discusses implications and limitations.
Section~\ref{sec:conclusion} concludes and suggests future work.

